% Appendix D.  Statement text follows the amended REV-PLAN contract.
% This is a section file: no document class, packages, or theorem declarations.
\section{Operator certificates for the local inequalities}\label{s:local-operator}

Fix a covered pattern \(\pi\in\mathcal P_{\rm ext}\) from
Definition~\ref{d:pattern}.  Write \(o\) for its out-neighbour and
\(f,g\) for the other two sites in
\(\{\mathsf e,\mathsf N,\mathsf W\}\), whether or not they are
removed.  For \(0<c<1\) set
\[
 F_c(m)=Z_{\rm cur}(\pi;m)-cZ_{\rm next}(\pi;m),
 \qquad m\in(\mathbb S^2)^{20}.
\]
All spin functionals in this appendix are finite sums with the laws
of Definition~\ref{d:moment}.  The proposition at the end gives a
sufficient certificate for this pattern's inequality in
Definition~\ref{h:H1}; it does not assert optimality of \(c\).
Empty products and tensor products mean \(1\) and \(\mathbb C\),
respectively, and the Pauli operator of a constant when there are
no vector variables is that scalar on \(\mathbb C\).

\subsection{Passage and leaf reductions}

\begin{lemma}\label{l:Iw}
Let \(Z_w\) be the finite template functional after deleting
the passage site and its incident factors, retaining all other
present bonds and all covered slot factors. Put \(P_{ab}=n_a\cdot n_b\)
for nonremoved sites and zero if either is removed.
For the out-neighbour \(o\) and the other two neighbours \(f,g\),
\[
 Z_{\rm cur}=Z_w[I],\qquad Z_{\rm next}=Z_w[J]>0,\qquad
 I=1+\tfrac13(P_{fg}+P_{fo}+P_{go}),\quad J=1+\tfrac15P_{fg}.
\]
For \(0<c<1\), \(F_c=Z_{\rm cur}-cZ_{\rm next}=Z_w[I-cJ]\), with
\[
 I-cJ=(1-c)+\tfrac13(P_{fo}+P_{go})+(\tfrac13-\tfrac c5)P_{fg}.
\]
Both partition polynomials are affine in each slot vector and
invariant under simultaneous orthogonal transformations of the slots.
\end{lemma}

\begin{proof}
For purposes of expanding a product, set \(n_x=0\) whenever \(x\)
is removed; this adds neither a variable nor a site law to the
functional.  By Definition~\ref{d:pattern}, the present passage
bonds in the current template go precisely to the nonremoved sites
among \(f,g,o\).  In the next template the passage law is
\(\lambda_{\mathrm D}\), its bond to \(o\) is deleted, and its
other present bonds are unchanged.  Consequently the factors
obtained by summing over the passage spin are, respectively,
\[
 \mathcal M\bigl((1+n\cdot n_f)(1+n\cdot n_g)
                         (1+n\cdot n_o)\bigr),\qquad
 \lambda_{\mathrm D}\bigl((1+n\cdot n_f)(1+n\cdot n_g)\bigr).
\]
The first and third moments vanish.  The second moment is
\(\delta_{ij}/3\) for \(\lambda_{\mathrm U}\) and
\(\delta_{ij}/5\) for \(\lambda_{\mathrm D}\), by
Lemma~\ref{l:moments}.  Expansion therefore gives exactly \(I\)
and \(J\), including when one or both of \(f,g\) are removed.
Finite summation can be performed in any order, proving both
partition identities.  Subtraction gives the asserted formula for
\(I-cJ\).

The network defining \(Z_w[1]\) has at most four incident factors
at each site, so \(Z_w[1]>0\) by Lemma~\ref{l:finitefacts}.
All its weights are nonnegative.  The remaining spins have norm
at most one, so \(P_{fg}\ge-1\), also with the removed-site
convention.  Thus
\[
 Z_{\rm next}=Z_w[J]\ge\frac45 Z_w[1]>0.
\]

Each slot occurs in just one linear factor.  All spin variables
in either original template occur in at most four factors:
the passage has at most three, each tip has its three slots and
at most one bond, each corner has its two slots and at most two
bonds, and each neighbour of the passage has at most four bonds.
For \(R\in\mathrm O(3)\), simultaneous transformation of slots
and spin variables leaves every scalar product unchanged.
The site-by-site polynomial covariance in
Lemma~\ref{l:moments}, with these separate degree bounds, then
shows that transforming the slots alone leaves each partition
polynomial unchanged.  This uses covariance of the indicated
polynomials, without asserting invariance of the finite laws on
arbitrary functions.
\end{proof}

\begin{lemma}\label{l:leaf}
In a covered template, integrating a nonremoved tip or corner
with no present internal bond gives a positive factor depending
only on its own slots.
A tip with its one internal bond to \(x\) gives
\(A(1+(3/5)u\cdot n_x)\), \(A>0,\ |u|\le1\).
A corner with exactly one internal bond to \(x\) gives
\(A(1+\rho u\cdot n_x)\), \(A>0,\ |u|\le1\), where
\(\rho=1/2\) for status \(\mathrm U\) and \(\rho=1/3\)
for status \(\mathrm D\).
\end{lemma}

\begin{proof}
All slot vectors in this proof are unit vectors.  A tip has law
\(\lambda_{\mathrm U}\), three slots, and at most one internal
bond, by Definition~\ref{d:pattern}.  If its slots are
\(a_1,a_2,a_3\), their normalizer is
\[
 A=\mathcal M\left(\prod_{j=1}^3(1+a_j\cdot n)\right)
   =1+\frac13\sum_{j<l}a_j\cdot a_l\ge\frac12,
\]
since \(0\le|a_1+a_2+a_3|^2=3+2\sum_{j<l}a_j\cdot a_l\).
With no bond this is the required positive factor.  With a bond
to \(x\), put
\[
 V=\mathcal M\left(n\prod_{j=1}^3(1+a_j\cdot n)\right).
\]
Summing over the tip gives \(A+V\cdot n_x\).  The normalized
mean bound in Lemma~\ref{l:35} gives \(|V|\le3A/5\), so
\(u=5V/(3A)\) has norm at most one and gives the claimed formula.
The moment in \(V\) has degree four, within the stated bound.

For a nonremoved corner let its two slots be \(a,b\), and put
\(t=a\cdot b\in[-1,1]\).  Write \(A\) for its slot normalizer
and \(V\) for its unnormalized slot mean.  Lemma~\ref{l:moments}
gives
\[
 \begin{array}{c|cc}
   &A&V\\ \hline
  \mathrm U&1+t/3&(a+b)/3\\
  \mathrm D&1+t/5&(a+b)/5 .
 \end{array}
\]
In particular \(A\ge2/3\) in status \(\mathrm U\), and
\(A\ge4/5\) in status \(\mathrm D\).  This proves the assertion
for corners with no internal bond.  If there is exactly one bond
to \(x\), its sum is again \(A+V\cdot n_x\).  As
\(|a+b|^2=2+2t\), in the two cases respectively we have
\[
 A^2-4|V|^2=\frac{(1-t)^2}{9}\ge0,
 \qquad
 A^2-9|V|^2=\frac{(1-t)(7-t)}{25}\ge0.
\]
Since \(A>0\), these inequalities give \(|V|\le A/2\) and
\(|V|\le A/3\).  Set \(u=V/(\rho A)\) with the corresponding
\(\rho\).  Every factor and vector constructed here depends only
on the slots at that leaf, as required.
\end{proof}

\begin{definition}\label{d:reduction}
Delete tips and corners with no present internal bond, dropping
their common positive factors. Replace every bonded tip by a
single vector and radius \(3/5\).
In the reduced representation also replace each one-bond corner
by a single vector of radius \(1/2\) or \(1/3\) according to
Lemma~\ref{l:leaf}; in the unreduced representation retain its spin
and two slots. Other bonded corners retain their two slots.
Let \(u=(u_1,\ldots,u_k)\) list these remaining vectors, and define
\(\widehat Z_{\rm cur},\widehat Z_{\rm next}\) by these replacements,
\(\widehat F_c=\widehat Z_{\rm cur}-c\widehat Z_{\rm next}\).
A group is one single vector or the two slots of a retained corner.
\end{definition}

Thus a single vector attached to a retained spin \(n_x\) means
the factor \(1+r_i u_i\cdot n_x\), where \(r_i\) is its prescribed
radius.  Each retained corner contributes
\((1+u_i\cdot n_z)(1+u_{i'}\cdot n_z)\) and retains its original
status law.  The functional after the passage sum consists of
these factors, the remaining internal bonds, and \(I\) or \(J\)
from Lemma~\ref{l:Iw}.  This specifies the polynomials for all
\(u_i\in\mathbb R^3\); no realizability condition on a replacement
vector is part of their definition.

\begin{corollary}\label{c:reduce}
In either representation, \(\widehat F_c\) is real, affine in
each vector and invariant under simultaneous \(\mathrm O(3)\).
If \(\widehat F_c\ge0\) on \((\mathbb S^2)^k\), then
\(F_c\ge0\) on \((\mathbb S^2)^{20}\).
For each original slot assignment there are a common \(K>0\)
and vectors \(\widehat u\) in the closed unit ball such that
\(Z_{\rm cur}=K\widehat Z_{\rm cur}(\widehat u)\) and
\(Z_{\rm next}=K\widehat Z_{\rm next}(\widehat u)\).
\end{corollary}

\begin{proof}
The slot factors and the replacement factors are real and affine
in each of their vector variables.  No vector is used twice.
Finite summation therefore preserves reality and separate
affineness.  In a term containing a pair insertion from \(I\) or
\(J\), each endpoint is a nonremoved non-arrival neighbour of the
passage.  Deleting its passage bond leaves at most three other
factors, and the insertion raises its degree by one.  Other
neighbours have degree at most three after passage deletion, and
retained corners have at most two slots and two bonds.  Replacing
a leaf changes a bond at its surviving endpoint to a single
linear factor and so does not increase that degree.  Hence all
summed spin variables have separate degree at most four.
Every factor is a scalar product or a constant with a real
radius.  Lemma~\ref{l:moments}, applied site by site as in
Lemma~\ref{l:Iw}, proves simultaneous \(\mathrm O(3)\) covariance
of both reduced partition polynomials, and thus of
\(\widehat F_c\).

Fix an original assignment \(m\in(\mathbb S^2)^{20}\).  No two
tips or corners are nearest neighbours in the template.  Thus
each sum in Lemma~\ref{l:leaf} can be performed with all other
spins fixed, without changing the slots of another leaf.  These
leaf factors contain no passage spin.  The same successive
summations therefore give, in the two partition functions, the
same product \(K>0\) of the leaf normalizers and the same list
\(\widehat u\) of replacement vectors and retained slots.
Lemma~\ref{l:leaf} bounds each replacement vector by one, and
the retained slots are unit vectors.  Consequently
\[
 Z_{\rm cur}=K\widehat Z_{\rm cur}(\widehat u),\qquad
 Z_{\rm next}=K\widehat Z_{\rm next}(\widehat u),\qquad
 F_c(m)=K\widehat F_c(\widehat u),\qquad |\widehat u_i|\le1.
\]
Removed sites have no factors, and an empty product of deleted
leaf normalizers is \(K=1\).

For completeness, if \(|v|\le1\), choose \(b=v/|v|\) when
\(v\ne0\), and choose \(b=(0,0,1)\) when \(v=0\).  Then
\[
 v=\frac{1+|v|}{2}b+\frac{1-|v|}{2}(-b).
\]
Successively using affineness in each vector expresses a value
of \(\widehat F_c\) on the product of balls as a sum of its
values on unit vectors with nonnegative product weights summing
to one.  For no vector variables this is the single value of the
constant polynomial.  Positivity on unit vectors therefore
implies positivity at \(\widehat u\), and the last displayed
identity proves the assertion.
\end{proof}

\subsection{Pauli operators and their symmetries}

We use the explicit matrices
\[
 \sigma_1=\begin{pmatrix}0&1\\1&0\end{pmatrix},\qquad
 \sigma_2=\begin{pmatrix}0&-\mathrm i\\\mathrm i&0\end{pmatrix},\qquad
 \sigma_3=\begin{pmatrix}1&0\\0&-1\end{pmatrix}.
\]
They are Hermitian, have square \(I\), and satisfy
\(\operatorname{Tr}\sigma_j=0\) and
\(\operatorname{Tr}(\sigma_j\sigma_l)=2\delta_{jl}\)
for \(1\le j,l\le3\).  Direct multiplication also gives
\(\sigma_j\sigma_l+\sigma_l\sigma_j=2\delta_{jl}I\).

\begin{definition}\label{d:Pauli}
For a real polynomial affine in each vector, expand uniquely
\[
 F(u)=\sum_{a\in\{0,1,2,3\}^k}f_a\prod_i u_i^{a_i},
 \qquad u_i^0=1,\qquad
 \mathcal O(F)=\sum_a f_a\bigotimes_i\sigma_{a_i},\quad\sigma_0=I.
\]
For \(\mathcal I\subset\{1,\ldots,k\}\), invert those vectors,
writing \(u^{\mathcal I}\) and \(F^{\mathcal I}(u)=F(u^{\mathcal I})\).
For \(|a|\le1\), put \(\rho(a)=(I+a\cdot\sigma)/2\).
Set \(S_z=\frac12\sum_i\sigma_3^{(i)}\) and
\(S_\pm=\sum_i(\sigma_1^{(i)}\pm\mathrm i\sigma_2^{(i)})/2\).
\end{definition}

Here \(u_i^j\) for \(j=1,2,3\) denotes a coordinate, not a power.
The expansion is unique because these are distinct monomials
in the \(3k\) real coordinates.  Separate affineness permits
at most one coordinate from each vector.  Its coefficients are
real, so \(\mathcal O(F)\) is Hermitian.  The map
\(F\mapsto\mathcal O(F)\) is real linear, commutes with finite
sums over auxiliary variables, and takes products on disjoint
sets of vector variables to tensor products of operators.

\begin{lemma}\label{l:inversion}
For \(|u_i|\le1\), each \(\rho(u_i)\) is positive semidefinite
of trace one, and
\[
 F(u)=\operatorname{Tr}\!\left[
      \mathcal O(F^{\mathcal I})
                 \bigotimes_i\rho((u^{\mathcal I})_i)\right].
\]
Thus a single inversion \(\mathcal I\) with
\(\mathcal O(F^{\mathcal I})\ge0\) implies \(F\ge0\)
on the product of closed unit balls, in particular on unit vectors.
\end{lemma}

\begin{proof}
The anticommutation identities give
\((a\cdot\sigma)^2=|a|^2I\).  For \(a\ne0\), the eigenvalues
of this traceless Hermitian matrix are \(|a|,-|a|\); for
\(a=0\) both are zero.  The eigenvalues of \(\rho(a)\) are
therefore \((1+|a|)/2,(1-|a|)/2\), which are nonnegative for
\(|a|\le1\).  Its trace is one.  The displayed trace identities
for the Pauli matrices give
\[
 \operatorname{Tr}(\sigma_j\rho(a))=a^j
       \quad(0\le j\le3),\qquad a^0=1.
\]
For any separately affine polynomial \(G\), multiplicativity
of the trace on tensor products now gives
\[
 \operatorname{Tr}\left[\mathcal O(G)
                          \bigotimes_i\rho(v_i)\right]
       =\sum_a g_a\prod_i v_i^{a_i}=G(v).
\]
Take \(G=F^{\mathcal I}\) and \(v=u^{\mathcal I}\).
Inverting the same coordinates twice is the identity, proving
the asserted formula.

The tensor product of the densities is positive semidefinite,
as follows by tensoring their orthonormal eigenbases.  If
\(\mathcal O(F^{\mathcal I})\) is positive semidefinite, choose
an orthonormal eigenbasis \((b_j)\) of it with eigenvalues
\(\lambda_j\ge0\).  For any positive semidefinite \(R\),
\[
 \operatorname{Tr}(\mathcal O(F^{\mathcal I})R)
       =\sum_j\lambda_j\langle b_j,Rb_j\rangle\ge0.
\]
Apply this with the product density.  The formulas also hold
for an empty tensor product, whose trace is one.
\end{proof}

\begin{lemma}\label{l:Ostruct}
For \(\mathcal I\) a union of vector groups and
\(\mathcal O=\mathcal O(\widehat F_c^{\mathcal I})\),
the matrix in the standard tensor basis is real symmetric.
It commutes with \(U^{\otimes k}\) for every \(U\in\mathrm{SU}(2)\)
and with \(S_z,S_+,S_-\).
For a retained \(\mathrm U\)-corner pair, if \(\Pi\) projects that
pair onto \(\operatorname{Sym}^2\mathbb C^2\), then
\(\mathcal O=\Pi\mathcal O\Pi\).
For a retained \(\mathrm D\)-corner pair, \(\mathcal O\) commutes
with exchange of its two factors.
\end{lemma}

\begin{proof}
Put \(G=\widehat F_c^{\mathcal I}\).  An inversion commutes with
a simultaneous orthogonal transformation, so Corollary~\ref{c:reduce}
implies \(G(Ru_1,\ldots,Ru_k)=G(u)\) for \(R\in\mathrm O(3)\).
Its coefficients are real and hence \(\mathcal O\) is Hermitian.
Entrywise conjugation fixes \(\sigma_0,\sigma_1,\sigma_3\) and
negates \(\sigma_2\).  Thus \(\overline{\mathcal O(G)}\) is the
Pauli operator of the polynomial obtained by reflecting the
second coordinate of every vector.  This is \(G\) itself by
its orthogonal covariance.  The operator is therefore real,
and Hermiticity makes it symmetric.

For \(U\in\mathrm{SU}(2)\), conjugation by \(U\) preserves the
real vector space of traceless Hermitian \(2\times2\) matrices
and its inner product \(\tfrac12\operatorname{Tr}(AB)\).
The Pauli matrices form an orthonormal basis for that space.
Consequently there is a real orthogonal matrix \(R\) with
\[
 U\sigma_jU^*=\sum_{l=1}^3R_{lj}\sigma_l.
\]
Extend it by \(R_{00}=1\), \(R_{0j}=R_{j0}=0\) for \(j>0\).
Expanding every tensor factor in the definition of the Pauli
operator gives
\[
 U^{\otimes k}\mathcal O(G)(U^*)^{\otimes k}
       =\mathcal O\bigl(G(R^{\mathsf T}u_1,\ldots,
                                      R^{\mathsf T}u_k)\bigr)
       =\mathcal O(G).
\]
Only orthogonality of \(R\) is needed here.

The concrete generator commutators follow by a polynomial
coefficient comparison.  For \(j=1,2,3\) and real \(t\), the
matrix
\[
 U_j(t)=\frac{I-\mathrm i t\sigma_j}{\sqrt{1+t^2}}
\]
is unitary of determinant one.  The just-proved commutation,
after multiplication by \((1+t^2)^{k/2}\), shows that
\(\mathcal O\) commutes with \((I-\mathrm i t\sigma_j)^{\otimes k}\)
for every real \(t\).  Its coefficient of \(t\) is
\(-\mathrm i\sum_i\sigma_j^{(i)}\), so \(\mathcal O\) commutes
with all three sums.  Their linear combinations give precisely
\(S_z,S_+,S_-\).  When \(k=0\) all these sums are zero.

Consider a retained corner pair \(i,i'\) at \(z\).
Since \(\mathcal I\) is a union of groups, its two slot factors
after inversion have the same sign \(\eta\in\{1,-1\}\):
\[
 (1+\eta u_i\cdot n_z)(1+\eta u_{i'}\cdot n_z).
\]
With all spins fixed, all other factors are independent of
\(u_i,u_{i'}\).  Real linearity of the Pauli map and its tensor
product property show that every summand of \(\mathcal O(G)\),
with the pair factors ordered first, has the form
\[
 (I+\eta n_z\cdot\sigma)\otimes
 (I+\eta n_z\cdot\sigma)\otimes H_n
\]
times a real scalar, for a matrix \(H_n\) on the other factors.
There is no positivity assumption on that scalar or on \(H_n\).
If \(z\) has status \(\mathrm U\), every value of \(n_z\) is
unit.  The matrix \(I+\eta n_z\cdot\sigma\) is then twice a
rank-one orthogonal projection, by the eigenvalue calculation
in Lemma~\ref{l:inversion}.  Its tensor square has range in
the symmetric square: if its one-dimensional range is spanned
by \(v\), the tensor square has range spanned by \(v\otimes v\).
It is Hermitian, so multiplication by the symmetric projection
on either side leaves it unchanged.  This proves
\(\Pi\mathcal O=\mathcal O=\mathcal O\Pi\), and hence
\(\mathcal O=\Pi\mathcal O\Pi\).

For status \(\mathrm D\), the value \(n_z=0\) is also allowed.
For every \(n_z\), including zero, exchanging the two factors
commutes with the tensor square of \(I+\eta n_z\cdot\sigma\).
The same holds after finite summation.  This proves the asserted
exchange commutation.  Each argument in this proof applies
separately to \(\widehat Z_{\rm cur}^{\mathcal I}\) and
\(\widehat Z_{\rm next}^{\mathcal I}\), as well as to their
linear combination: both polynomials have the covariance and
the corner slot factors used above.
\end{proof}

\subsection{Exchange sectors and the central weight}

\begin{definition}\label{d:exchange}
Let \(V\) be the tensor product of \(\mathbb C^2\) for singles,
\(\operatorname{Sym}^2\mathbb C^2\) for retained \(\mathrm U\)-pairs,
and \(\mathbb C^2\otimes\mathbb C^2\) for retained \(\mathrm D\)-pairs.
For each assignment \(\epsilon\) of exchange signs to the
\(\mathrm D\)-pairs, let \(V_\epsilon\) replace each of those
pair factors by its symmetric or exterior square.
If there are \(t\) singles, let \(\kappa=0\) for even \(t\),
\(\kappa=1/2\) for odd \(t\), and
\(V_\epsilon^\circ=\ker(S_z-\kappa I)|_{V_\epsilon}\).
The standard sector basis uses \(e_\pm\) for singles and
\(e_+\otimes e_+\), \(e_+\otimes e_-+e_-\otimes e_+\),
\(e_-\otimes e_-\) for symmetric pairs, with
\(e_+\otimes e_--e_-\otimes e_+\) for antisymmetric pairs.
\end{definition}

We realize all these spaces as subspaces of
\((\mathbb C^2)^{\otimes k}\) with its usual Hermitian inner
product; the pair basis vectors are deliberately unnormalized.
The symmetric and exterior squares are orthogonal, so
\(V=\bigoplus_\epsilon V_\epsilon\) is an orthogonal direct sum.
The central-sector basis is the subset of the indicated tensor
basis having total weight \(\kappa\).  These basis vectors are
orthogonal, have positive squared norms that are powers of two,
and have coordinates in \(\{0,1,-1\}\) in the full tensor basis.

The three total spin operators preserve each such tensor
sector.  Indeed, on a pair put
\[
 v_+=e_+\otimes e_+,\quad
 v_0=e_+\otimes e_-+e_-\otimes e_+,\quad
 v_-=e_-\otimes e_-,\quad
 w=e_+\otimes e_--e_-\otimes e_+.
\]
The pair's \(S_z\)-weights on \(v_+,v_0,v_-,w\) are
\(1,0,-1,0\), respectively, and its nonzero ladder actions are
\[
 S_-v_+=v_0,\quad S_-v_0=2v_-,\qquad
 S_+v_-=v_0,\quad S_+v_0=2v_+.
\]
Both ladders kill \(w\); all ladder actions not displayed on
these four vectors are zero.  On a single factor the weights
are \(1/2,-1/2\), and \(S_-e_+=e_-\), \(S_+e_-=e_+\),
with the other two ladder actions zero.  The total operators
are sums of these local actions.  In particular every weight
in a tensor sector is in \(\kappa+\mathbb Z\).  On each sector,
direct multiplication of the single-factor matrices gives
\[
 S_z^*=S_z,\qquad S_+^*=S_-,\qquad
 [S_z,S_\pm]=\pm S_\pm,\qquad [S_+,S_-]=2S_z.
\]
Operators on different tensor factors commute, so the same
identities hold for their sums and their sector restrictions.

\begin{lemma}\label{l:sector}
Let \(V'\) be any tensor sector formed from singles and
symmetric or exterior pair factors as above.
If a Hermitian \(\mathcal O'\) on \(V'\) commutes with
\(S_z,S_+,S_-\), then it preserves the central-weight space
\(V'^\circ\), and its restriction there has the same set of
eigenvalues as \(\mathcal O'\).
In particular \(\mathcal O'\) is positive definite if and only if
that restriction is positive definite.
\end{lemma}

\begin{proof}
Commutation with \(S_z\) immediately implies preservation of
\(V'^\circ\).  Every eigenvalue of the restriction is therefore
an eigenvalue on \(V'\).  Conversely, let \(\lambda\) be an
eigenvalue of \(\mathcal O'\), and let its nonzero eigenspace
be \(E_\lambda\).  All three spin operators preserve
\(E_\lambda\).  The restriction of \(S_z\) to this space is
Hermitian, so there is a nonzero \(v\in E_\lambda\) with
\(S_zv=mv\).  The explicit tensor weights above show that
\(m\in\kappa+\mathbb Z\).

For any vector of weight \(m\), the adjoint and commutator
identities give
\[
 \|S_-v\|^2-\|S_+v\|^2
       =\langle v,[S_+,S_-]v\rangle=2m\|v\|^2.
\]
Thus \(S_-v\ne0\) if \(m>0\), and \(S_+v\ne0\) if \(m<0\).
The commutators with \(S_z\) show that these vectors have weights
\(m-1\) and \(m+1\), respectively.  Commutation with
\(\mathcal O'\) keeps every nonzero ladder vector in
\(E_\lambda\).

If \(m>\kappa\), then \(m-\kappa\) is a positive integer.
Apply \(S_-\) that many times.  Every weight before the final
step is strictly positive, so no step vanishes, and the final
vector has weight \(\kappa\).  If \(m<\kappa\), then
\(\kappa-m\) is a positive integer.  Every weight before the
final step of successive raising is negative: it is at most
\(-1\) for \(\kappa=0\), and at most \(-1/2\) for
\(\kappa=1/2\).  Hence raising \(\kappa-m\) times also gives a
nonzero vector of weight \(\kappa\).  If \(m=\kappa\), take
\(v\) itself.  In all cases \(\lambda\) occurs on
\(V'^\circ\).

The two Hermitian operators consequently have the same set of
eigenvalues.  A finite-dimensional Hermitian operator is
positive definite exactly when all its eigenvalues are positive,
by its orthonormal eigenbasis, proving the last assertion.
This argument asserts nothing about eigenvalue multiplicities.
\end{proof}

\subsection{Integer positivity and reconstruction}

\begin{lemma}\label{l:gram-residual}
Let \(M\) be a symmetric integer matrix of size \(n\).
Let \(g,d,S\) be positive integers and \(L\) an integer
lower-triangular \(n\times n\) matrix; no invertibility is required.
Define the symmetric integer residual
\[
 E=dS^2M-gLL^{\mathsf T}.
\]
If, for every row,
\[
 E_{ii}>0,\qquad E_{ii}>\sum_{j\ne i}|E_{ij}|,
\]
then \(E\) and \(M\) are positive definite over \(\mathbb R\)
and, as Hermitian matrices, over \(\mathbb C\).
The case \(L=0\), \(g=d=S=1\), is the direct strict
diagonal-dominance certificate.
All hypotheses can be checked with integer sums, products,
absolute values and comparisons, without division or determinants.
\end{lemma}

\begin{proof}
Let \(z\in\mathbb C^n\).  Symmetry and reality of \(E\) imply
\begin{align*}
 z^*Ez
 &=\sum_iE_{ii}|z_i|^2+
           2\sum_{i<j}E_{ij}\operatorname{Re}(\overline z_i z_j)\\
 &\ge\sum_iE_{ii}|z_i|^2-
           2\sum_{i<j}|E_{ij}|\,|z_i|\,|z_j|\\
 &\ge\sum_i\left(E_{ii}-\sum_{j\ne i}|E_{ij}|\right)|z_i|^2.
\end{align*}
The last step uses \(2ab\le a^2+b^2\) for nonnegative \(a,b\).
Every parenthesized coefficient is positive, so \(z^*Ez>0\)
for every nonzero \(z\).  Also
\[
 dS^2 z^*Mz=z^*Ez+g z^*LL^{\mathsf T}z
             =z^*Ez+g\|L^{\mathsf T}z\|^2>0
             \qquad(z\ne0).
\]
Since \(dS^2>0\), this proves positivity of \(M\) over
\(\mathbb C\), and restriction proves it over \(\mathbb R\).
The argument also covers size zero, when positive definiteness
is vacuous.  No rank condition on \(L\) was used.

In particular, lower triangularity lets one form its Gram entries
as the integer sums
\[
 (LL^{\mathsf T})_{ij}
       =\sum_{a=1}^{\min(i,j)}L_{ia}L_{ja}
       \quad(1\le i,j\le n).
\]
The hypotheses and residual entries therefore require only the
listed integer operations and comparisons.  Setting \(L=0\)
and \(g=d=S=1\) leaves \(E=M\), giving the direct certificate.
\end{proof}

\begin{lemma}\label{l:integercheck}
For \(c=p/q\in(0,1)\) rational with positive integers \(p,q\),
choose a positive integer \(D_0\) clearing the coefficients of
\(3D_0\widehat Z_{\rm cur}\) and \(5D_0\widehat Z_{\rm next}\).
In the integer-coordinate central-sector basis, let \(A_\epsilon\)
and \(A'_\epsilon\) be the Hermitian-form matrices of their
inverted Pauli operators. These are symmetric integer matrices.
The form of \(\mathcal O(\widehat F_c^{\mathcal I})\) on the sector
is positive definite if and only if
\[
 M_\epsilon=5qA_\epsilon-3pA'_\epsilon
\]
is positive definite. The primary verification uses the integer
Gram/residual witness of Lemma~\ref{l:gram-residual}.
As an optional fallback only, positivity is equivalent to positivity of all its
leading principal minors. More explicitly, for a size-\(t\) matrix
set \(A^{(0)}=M_\epsilon\), \(p_{-1}=1\), and for \(0\le j<t\)
set \(p_j=A^{(j)}_{jj}\) and update the trailing entries by
\[
 A^{(j+1)}_{il}
    =\frac{p_jA^{(j)}_{il}-A^{(j)}_{ij}A^{(j)}_{jl}}{p_{j-1}}
       \quad(i,l>j).
\]
As long as preceding pivots are nonzero, each division is exact
in the integers and \(p_j\) is the leading principal minor of
order \(j+1\). Acceptance requires every division exact and every
pivot positive.
\end{lemma}

\begin{proof}
We first specify a reconstruction directly from the pattern;
this also gives an explicit choice of denominator.  Let \(X\)
be the set of spins still summed after the passage and the
chosen leaf reductions, and let \(B\) be their remaining present
internal bonds.  For each remaining vector \(u_i\), write
\(x(i)\in X\) for the spin to which it is attached.  Put
\(r_i=1\) for a retained corner slot and give a single vector
its radius \(r_i\in\{3/5,1/2,1/3\}\) from
Definition~\ref{d:reduction}.  The common remaining weight is
\[
 W(n,u)=\prod_{\{x,y\}\in B}(1+n_x\cdot n_y)
                     \prod_{i=1}^k(1+r_i u_i\cdot n_{x(i)}).
\]
If \(\mathcal T\) denotes the product of the original status
laws on \(X\), Lemma~\ref{l:Iw} and the definition of the
representation give
\[
 \widehat Z_{\rm cur}(u)=\mathcal T[W(n,u)I(n)],\qquad
 \widehat Z_{\rm next}(u)=\mathcal T[W(n,u)J(n)].
\]
All nonremoved sites among \(f,g,o\) remain in \(X\); they are
never periphery leaves.  Isolated spins may be retained in
\(X\), contributing their zeroth moment one.

For a word \(w=(j_1,\ldots,j_s)\) of coordinate indices in
\(\{1,2,3\}\), define the local moment
\(\mu_{\mathrm{st}}(w)=\lambda_{\mathrm{st}}
(n_{j_1}\cdots n_{j_s})\).  Only words of length at most four
will occur.  Their exact values, by Lemma~\ref{l:moments}, are
\[
 \begin{array}{c|cc}
 w&\mu_{\mathrm U}(w)&\mu_{\mathrm D}(w)\\ \hline
 \varnothing&1&1\\
 (i)\text{ or }(i,j,l)&0&0\\
 (i,j)&\delta_{ij}/3&\delta_{ij}/5\\
 (i,j,l,m)&
 \dfrac{\delta_{ij}\delta_{lm}+\delta_{il}\delta_{jm}
                         +\delta_{im}\delta_{jl}}{15}&
 \dfrac{\delta_{ij}\delta_{lm}+\delta_{il}\delta_{jm}
                         +\delta_{im}\delta_{jl}}{25}.
 \end{array}
\]
In particular every \(\mathrm U\) moment has denominator dividing
\(15\), and every \(\mathrm D\) moment has denominator dividing
\(25\).

Here is an entrywise formula for the entire coefficient tensor.
For \(\nu\in\{\mathrm{cur},\mathrm{next}\}\), let
\(\mathcal H_\nu\) contain an empty insertion of weight one.
For \(\nu=\mathrm{cur}\), also include \((v,w,j)\) of weight
\(1/3\) for each pair \(vw\in\{fg,fo,go\}\) whose endpoints are
nonremoved and each \(1\le j\le3\).  For
\(\nu=\mathrm{next}\), include just the nonremoved pair \(fg\)
with these three coordinate choices, each of weight \(1/5\).
Denote the insertion weight by \(\omega_\nu(h)\).
Fix a vector multi-index \(a\in\{0,1,2,3\}^k\) and bond
indices \(b\in\{0,1,2,3\}^{B}\).  At \(x\in X\), form a word
\(w_x(a,b,h)\) by collecting:
the index \(a_i\) for each \(i\) with \(x(i)=x\) and
\(a_i\ne0\); the index \(b_e\) for each incident bond \(e\)
with \(b_e\ne0\); and the insertion coordinate \(j\) if
\(h=(v,w,j)\) and \(x\) is one of \(v,w\).
The order of these indices is immaterial.  Expansion of each
factor as a sum of its constant and three coordinate terms gives
\[
 z_{\nu,a}=
 \left(\prod_{i=1}^k r_i^{\mathbf1_{\{a_i\ne0\}}}\right)
 \sum_{b\in\{0,1,2,3\}^{B}}\ \sum_{h\in\mathcal H_\nu}
 \omega_\nu(h)\prod_{x\in X}
          \mu_{\mathrm{st}(x)}\bigl(w_x(a,b,h)\bigr),
 \qquad
 \widehat Z_\nu(u)=\sum_a z_{\nu,a}\prod_i u_i^{a_i}.
\]
Indeed, each choice of terms from the product determines just
such \(a,b,h\); conversely each such triple determines one term,
with the indicated scalar radius and insertion weights.  The
product functional of that term is the product of its local
moments.  This proves the formula by finite distributivity.
The degree check in Corollary~\ref{c:reduce} shows explicitly
that every \(w_x\) has length at most four.

Write \(r_i=\alpha_i/\beta_i\) with positive integers
\(\alpha_i,\beta_i\), using \(\beta_i=1\) for retained slots.
If \(u\) and \(v\) are the numbers of \(\mathrm U\) and
\(\mathrm D\) sites in \(X\), a permissible positive denominator
is
\[
 D_0=15^u25^v\prod_{i=1}^k\beta_i.
\]
Each summand in the coefficient formula, after multiplication
by \(3D_0\) for the current polynomial or \(5D_0\) for the next,
is an integer: a vector radius is used at most once, one local
moment is used at each site, and the only remaining insertion
denominator is \(3\) or \(5\), respectively.  Any smaller
positive \(D_0\) is equally valid if the two coefficient
integrality conditions are verified.  Thus the existence and
positivity of a denominator follow from the pattern rather than
from stored matrix data.

For clarity, the remaining Pauli and basis transformations also
have exact entrywise formulas.  Set
\[
 T_{{\rm cur},a}=3D_0z_{{\rm cur},a},\qquad
 T_{{\rm next},a}=5D_0z_{{\rm next},a},\qquad
 \chi_{\mathcal I}(a)=(-1)^{|\{i\in\mathcal I:a_i\ne0\}|}.
\]
Index the full tensor basis by
\(s,t\in\{+,-\}^k\).  For either polynomial the inverted,
scaled Pauli operator has entries
\[
 (P_\nu)_{st}=\sum_a
       T_{\nu,a}\chi_{\mathcal I}(a)
                         \prod_{i=1}^k(\sigma_{a_i})_{s_i t_i}.
\]
Let \(B_\epsilon\) be the rectangular matrix whose columns are
the integer-coordinate basis vectors of \(V_\epsilon^\circ\)
in Definition~\ref{d:exchange}, expanded in this full tensor
basis, in any fixed order.  The requested matrices are exactly
\[
 A_\epsilon=B_\epsilon^*P_{\rm cur}B_\epsilon,
 \qquad A'_\epsilon=B_\epsilon^*P_{\rm next}B_\epsilon.
\]
The moment formula, inversion signs, explicit Pauli entries, and
explicit basis columns therefore reconstruct every matrix entry
by finite sums and products from the pattern and the chosen
representation.  No matrix or pivot list is an additional
mathematical input.

Since the \(T_{\nu,a}\) are integers and the Pauli matrices
have entries in \(\mathbb Z[\mathrm i]\), the \(P_\nu\) have
entries in \(\mathbb Z[\mathrm i]\).  The reality and symmetry
proof of Lemma~\ref{l:Ostruct} applies separately to each
polynomial, so the \(P_\nu\) are real symmetric.  Their entries
are thus integers.  The columns of \(B_\epsilon\) are real
integer vectors, proving that \(A_\epsilon,A'_\epsilon\) are
symmetric integer matrices.  Coefficient reflection in the
second coordinate alternatively shows that only terms with
an even number of \(\sigma_2\) factors can be nonzero, so the
same construction can be evaluated with real integer entries.

Linearity of the Pauli map gives the exact identity
\begin{align*}
 M_\epsilon
 &=5qA_\epsilon-3pA'_\epsilon\\
 &=15D_0q\,B_\epsilon^*
                \mathcal O(\widehat F_c^{\mathcal I})B_\epsilon.
\end{align*}
The columns of \(B_\epsilon\) form a basis of the central sector,
so \(B_\epsilon\) is an isomorphism from its coordinate space
onto that sector.  The last matrix on the right is the matrix
of the restricted Hermitian form, even though the basis is
unnormalized.  Since \(15D_0q>0\), its form is positive definite
exactly when \(M_\epsilon\) is positive definite.  This is the
matrix to which Lemma~\ref{l:gram-residual} applies in the primary
verification.

We finish by proving the optional determinant test and its
fraction-free implementation.  For any real symmetric matrix,
positive definiteness is equivalent to positivity of every
leading principal minor.  One elementary proof is induction on
the size.  If the first diagonal entry is \(a>0\), write the
matrix as
\[
 M=\begin{pmatrix}a&r^{\mathsf T}\\r&C\end{pmatrix},
 \qquad Q=C-rr^{\mathsf T}/a.
\]
Completing the square gives
\[
 (x,y)^{\mathsf T}M(x,y)
       =a(x+r^{\mathsf T}y/a)^2+y^{\mathsf T}Qy.
\]
Thus \(M\) is positive definite exactly when \(a>0\) and \(Q\)
is positive definite.  A block row elimination of determinant
one shows that every leading minor of size \(r\ge2\) of \(M\)
is \(a\) times the leading minor of size \(r-1\) of \(Q\).
Induction proves both directions, including size one; size zero
has no conditions.  This also proves the Hermitian-form version
for a real matrix, since for \(z=x+\mathrm i y\) its Hermitian
quadratic form is \(x^{\mathsf T}Mx+y^{\mathsf T}My\).

For the recurrence in the statement, index rows and columns
by \(0,\ldots,t-1\), and let \(\Delta_j\) be the determinant
of the leading \(j\times j\) block of \(M_\epsilon\), with
\(\Delta_0=1\).  For \(i,l\ge j\) the induction invariant is
\[
 A^{(j)}_{il}=
  \det M_\epsilon[(0,\ldots,j-1,i),(0,\ldots,j-1,l)].
\]
The ordered lists here specify rows and columns in that order;
for \(j=0\) this is the original entry.  In particular
\(p_j=\Delta_{j+1}\) and \(p_{j-1}=\Delta_j\).
Assume the preceding pivots are nonzero.  The bordered-minor
identity is
\begin{align*}
 &A^{(j)}_{jj}A^{(j)}_{il}
                    -A^{(j)}_{ij}A^{(j)}_{jl}\\
 &\hspace{1cm}=\Delta_j\,
  \det M_\epsilon[(0,\ldots,j,i),(0,\ldots,j,l)]
                       \qquad(i,l>j).
\end{align*}
For \(j=0\) it is the \(2\times2\) determinant formula.  For
\(j>0\), the leading block \(B_j\) of size \(j\) is invertible
because \(\det B_j=\Delta_j\ne0\).  In the submatrix with the
two added row and column indices, eliminate the entries below
\(B_j\).  If the resulting \(2\times2\) Schur block has entries
\(q_{ab}\), each bordered \((j+1)\)-minor is
\(\Delta_j q_{ab}\), while the full \((j+2)\)-minor is
\(\Delta_j(q_{jj}q_{il}-q_{ij}q_{jl})\).
Substitution proves the displayed identity.  The quotient by
\(p_{j-1}=\Delta_j\) is therefore the indicated next bordered
minor, an integer.  This proves exact divisibility, the induction
invariant, and the pivot formula throughout the recurrence.
All positive pivots consequently imply all positive leading
minors, and conversely a positive definite matrix yields exact
divisions and positive pivots at every step.  These determinant
calculations are needed only for this optional fallback.
\end{proof}

\begin{proposition}\label{p:certificate}
Let \(c=p/q\in(0,1)\), with \(p,q\in\mathbb Z_{>0}\).
Choose either allowed representation, an inversion which is a
union of groups, and a justified positive denominator \(D_0\).
For every supported central exchange sector reconstruct
\(M_\epsilon=5qA_\epsilon-3pA'_\epsilon\) as in
Lemma~\ref{l:integercheck}.
Suppose that for each sector there are positive integers
\(g_\epsilon,d_\epsilon,S_\epsilon\) and an integer lower-triangular
matrix \(L_\epsilon\) of its dimension such that
\[
 E_\epsilon=d_\epsilon S_\epsilon^2M_\epsilon
                  -g_\epsilon L_\epsilon L_\epsilon^{\mathsf T},
 \qquad
 (E_\epsilon)_{ii}>
          \sum_{j\ne i}|(E_\epsilon)_{ij}|\quad\text{for every }i.
\]
Then \(F_c(m)\ge0\) for every \(m\in(\mathbb S^2)^{20}\).
The operator may vanish on the discarded antisymmetric
\(\mathrm U\)-pair subspaces; positive definiteness is required
only on the supported exchange sectors.
\end{proposition}

\begin{proof}
Let \(\mathcal O=\mathcal O(\widehat F_c^{\mathcal I})\).
For every sector the residual is real symmetric, since its
reconstructed \(M_\epsilon\) and its Gram matrix are symmetric.
Its strict row inequality forces
\((E_\epsilon)_{ii}>0\), because the sum of absolute values is
nonnegative.  All hypotheses of Lemma~\ref{l:gram-residual}
therefore hold, and every \(M_\epsilon\) is positive definite.
Lemma~\ref{l:integercheck}, including its positive scale
\(15D_0q\), makes the form of \(\mathcal O\) positive definite
on every \(V_\epsilon^\circ\).

Let \(\Pi_V\) be the product of the symmetric projections for
all retained \(\mathrm U\)-pairs, tensored with the identities
on all remaining groups.  These projections act on disjoint
factors and commute.  Lemma~\ref{l:Ostruct} gives
\(\mathcal O=\Pi_V\mathcal O\Pi_V\), with
\(\operatorname{Ran}\Pi_V=V\).  In particular \(\mathcal O\)
vanishes on \(V^\perp\) and maps \(V\) into itself.
Its commutation with every retained \(\mathrm D\)-pair exchange
also makes every \(V_\epsilon\) invariant.  The concrete spin
operators preserve these sectors, as checked after
Definition~\ref{d:exchange}; Lemma~\ref{l:Ostruct} shows that
the restriction of \(\mathcal O\) to each sector commutes with
their restrictions.  Lemma~\ref{l:sector} now promotes positive
definiteness on \(V_\epsilon^\circ\) to positive definiteness
on \(V_\epsilon\).

The direct sum \(V=\bigoplus_\epsilon V_\epsilon\) is orthogonal.
For any full tensor vector \(z\), write
\(\Pi_Vz=\sum_\epsilon z_\epsilon\) with
\(z_\epsilon\in V_\epsilon\).  Invariance and orthogonality give
\[
 \langle z,\mathcal Oz\rangle
       =\sum_\epsilon\langle z_\epsilon,
                                  \mathcal Oz_\epsilon\rangle\ge0.
\]
Thus \(\mathcal O\) is positive semidefinite on the full tensor
space, allowing precisely the discarded antisymmetric
\(\mathrm U\)-pair subspaces to have zero action.
Lemma~\ref{l:inversion} implies \(\widehat F_c\ge0\) on the
product of closed unit balls, hence on unit vectors.
Corollary~\ref{c:reduce} gives \(F_c(m)\ge0\) for every original
twenty-slot assignment, proving the proposition.
\end{proof}

